% =====================================================================
%  GMN2 §8 不规则奇点 + §9.1–9.3 线性 SU(2) 箭图的标度极限
%  原文：wkb-paper.tex 5955–6356 行；印刷页 p.100–106（PDF p.101–107）
%  编号公式 (8.1)–(8.5)、(9.1)–(9.15) 共 20 条；插图 2 幅（图 42、图 43）
%  备注：原文 §9.3 末尾的 align 在 PDF 上占两个编号 (9.14)(9.15)，
%        故本块编号式为 5 + 15 = 20 条（BRIEF 表估 18 条，实测 20 条）。
%        跨块引用一律硬编码原文编号并加 % TODO；块内引用用 \eqref。
% =====================================================================

\section{不规则奇点}\label{sec:irregular}

过去三节我们一直关注正则奇点的情形，但这些构造也可以搬到不规则奇点的情形。故事与正则情形颇为相似，只需作如下修改：

\begin{enumerate}

\item 我们首先给出包含不规则奇点时的三角剖分的适当概念。设 $\CP_*$ 是一个不规则奇点，$\lambda^2$ 在那里有 $L+2$ 阶极点，其中 $L \geq 1$ 为整数。（第 3 节% TODO 合稿改 \ref{sec:physics}
对秩为二的 Hitchin 系统给出的边界条件涉及 $L = 1$ 的不规则奇点；而在第 9 节% TODO 合稿改 \ref{section:Superconformal}
中我们会遇到 $L = N+2$、$N\geq 0$ 的奇点。）我们在 $\CP_*$ 周围画一个圆 $S^1(\CP_*)$，它围出一个圆盘 $D(\CP_*)$，视之为无穷小。在这个圆上标出 $L$ 个点 $Q_i$，$i=1, \dots, L$。这些点的循环次序规定为：$\dots, Q_i, Q_{i+1}, \dots$ 绕 $\CP_*$ 顺时针排列。于是我们的三角剖分其实是曲面 $C'$ 的三角剖分，而 $C'$ 由把所有圆盘 $D(\CP_*)$ 挖掉得到。顶点是所有不规则奇点周围的标号点 $Q_i$，再加上全部正则奇点 $\CP_i$。边必须包含圆周 $S^1(\CP_*)$ 上连接相邻点 $Q_i$ 的那些弧段。我们把这些弧段称为\emph{边界边}；它们在下面会有特殊地位。

\item 为了定义一个\emph{装饰}三角剖分，我们需要在每个顶点附近为联络 $\CA$ 选一个平坦截面（只定到标度因子）。对不规则奇点，这意味着在 $S^1(\CP_*)$ 上的每个点 $Q_i$ 附近各选一个平坦截面。与正则情形一样，我们希望把这一选择缩小为离散的选取。为此我们注意到：在不规则奇点的情形，除可能存在的单值之外，还有斯托克斯现象。\footnote{以下我们假定：复曲线上\emph{亚纯}联络的标准斯托克斯理论可以按最明显的方式推广，以适用于联络 $\CA$——它平坦，但并不亚纯。关于这一确切情形，我们没有找到任何文献，不过 \cite{MR1864833} 中有一些多少相关的构造。}从 $\CP_*$ 发出 $L$ 条射线（「斯托克斯射线」），围出张角为 $2\pi/L$ 的扇区。一个平坦截面若沿着斯托克斯射线某一侧的、指向 $\CP_*$ 的射线当 $z \to \CP_*$ 时渐近地指数式小，那么它在斯托克斯射线的另一侧就变成指数式大。\footnote{在关于斯托克斯现象的文献中有两类射线，分别称为 \emph{Stokes} 与 \emph{anti-Stokes}，二者在理论的系统展开中都起重要作用。遗憾的是，这一术语在各家作者的用法中并不一致。按我们的约定，标准 Airy 函数 $Ai(x)$ 沿正实轴有实的指数衰减，沿负实轴则有多项式型衰减并带振荡包络；它的斯托克斯射线位于负实轴以及 $\vert {\rm arg}(x)\vert = \pi/3$ 处。}我们这样定义一个装饰：取 $Q_i$ 附近的一个平坦截面（模去整体重标度），使得绕 $\CP_*$ 作解析延拓之后，它当 $z \to \CP_*$ 时在由斯托克斯射线围成的\emph{某个}扇区内指数式小。

\item 这样的截面（模去重标度）可以排成一个无穷序列
\begin{equation}\label{eq:bisequenceStokes}
\cdots, s_{-3}, s_{-2}, s_{-1}, s_0, s_1, s_2, s_3, \dots
\end{equation}
其排序规定为：若 $s_n$ 是扇区 $\CS$ 中的小解，则 $s_{n+1}$ 就是沿顺时针方向的下一个扇区中的小解。（若单值平凡，则可选取 $s_n$ 的标度，使该序列以 $L$ 为周期。）此外我们进一步限制装饰的选取：若在点 $Q_i$ 处选了 $s_n$，那么在点 $Q_{i+1}$ 处也必须选 $s_{n+1}$，如此递推。于是，不规则奇点处装饰的选取归结为在\emph{一个}标号点上选一个平坦截面，而不是在每个点上独立地选。因此，不规则奇点附近一切可能的装饰构成一个 $\IZ$-挠子；若单值平凡，它可以化归为 $\IZ_L$-挠子。

\item Fock–Goncharov 坐标 $\CX_E^T$ 的定义现在可以照旧给出，只需注意一条重要约定：若 $E$ 是边界边，则规定 $\CX_E^T = 0$。

\item 不规则奇点附近 WKB 叶状结构的局部行为，与正则奇点附近相当不同。每条 WKB 曲线不再各向同性地盘旋进入奇点，而是渐近地切于 $L$ 条射线之一。这些射线（相位为 $\vartheta$）我们称为 \emph{WKB 射线}，它们在不规则奇点 $\CP_*$ 周围的无穷小圆 $S^1(\CP_*)$ 上决定 $L$ 个点 $Q_i$，$i=1, \dots, L$；这些正是我们定义 $T_\WKB$ 时要用的标号点。若 $\vartheta = \arg \zeta$，则相位为 $\vartheta$ 的 WKB 射线与联络 $\CA$ 的反斯托克斯射线重合。
\insfig{2.7in}{petals}{$L = 4$ 的不规则奇点附近 WKB 叶状结构的局部行为。WKB 曲线聚集到这 $4$ 条 WKB 射线上，相邻射线之间相隔 $\pi/2$ 弧度的弧段。在奇点的小邻域内，WKB 曲线形如花瓣，把相邻的 WKB 射线连接起来。}

\item $T_\WKB$ 的定义本质上与之前一样。若一条 WKB 曲线渐近地趋向某条终止于不规则奇点 $\CP_*$ 的 WKB 射线，我们就认为它终止于 $S^1(\CP_*)$ 上相应的点 $Q_i$。于是照常，分界 WKB 曲线把 $C'$ 切成若干胞腔，每个胞腔由一般的 WKB 曲线叶状铺满，而 $T_\WKB$ 的边由每个胞腔中取一条一般 WKB 曲线构成。注意有些一般 WKB 曲线完全落在不规则奇点的任意小邻域内，并把相邻的 WKB 射线连接起来，如图 42 所示% TODO 合稿改 \ref{fig:petals}
。$T_\WKB$ 的边中有 $L$ 条这样的曲线，它们就是上面提到的、连接相邻 $Q_i$ 的边界边。每条这样的边围出一个触及 $\CP_*$ 的花瓣状区域；圆盘 $D(\CP_*)$ 就是这 $L$ 个花瓣的并。

\item $T_\WKB$ 在顶点 $Q_j$ 处的装饰，取这样的平坦截面：沿穿过 $Q_j$、指向 $\CP_*$ 的 WKB 射线走时，它变得指数式小。让我们把这个截面描述得更精确些。把 $\CP_*$ 放在 $z = \infty$，并设那里 $\lambda^2 \sim z^{L-2} dz^2$。$L$ 条 WKB 射线位于
\begin{equation}\label{eq:PolyRays}
r_j := \left\{ \arg(z) = \frac{2}{L} \vartheta + \frac{2\pi j}{L} \right\}, \quad j=1, \dots, L.
\end{equation}
当 $z \to \infty$ 时，平坦截面的形式渐近式形如
\begin{equation}\label{eq:ch11-formal-asymp}
\exp \left[ \pm \frac{2}{L}\frac{R}{\zeta} z^{\half L} + \cdots \right] s_{const}.
\end{equation}
特别地，其范数由指数实部的符号控制，而该符号在斯托克斯射线
\begin{equation}\label{eq:ch11-stokes-rays}
{\rm arg}(z) = \frac{2}{L}\vartheta_\zeta + \frac{(2j+1)\pi}{L}, \qquad j = 0, \dots, L-1
\end{equation}
两侧发生改变，其中我们记 $\vartheta_\zeta = \arg \zeta$。只要 $e^{i\vartheta} \not= - e^{i\vartheta_\zeta}$，由这些斯托克斯射线围成的每个扇区中都恰含一条 WKB 射线 $r_j$。沿这条 WKB 射线，一般平坦截面的范数渐近于
\begin{equation} \label{eq:norm-sign}
\exp\left[ \pm \frac{ 2 R}{L} \frac{\vert z\vert^{L/2}}{\vert \zeta\vert} \Re \left(e^{i (\vartheta + \pi j - \vartheta_\zeta)} \right) + \cdots \right].
\end{equation}
若在 \eqref{eq:norm-sign} 中取的 $\pm$ 号与 $\Re\left(e^{i (\vartheta + \pi j - \vartheta_\zeta)} \right)$ 的符号相反，则对任何 $\zeta \in \IH_\vartheta$，该范数沿 $r_j$ 当 $z \to \infty$ 时指数式小。到此为止我们只讨论了形式渐近；但有一条重要的原理：存在唯一的平坦截面（称为「小截面」），它的范数沿 $r_j$ 确实具有这种指数式小的渐近行为（对 $\zeta \in \IH_\vartheta$）。我们就选这个小截面作为 $T_\WKB(\vartheta)$ 在 WKB 射线 $r_j$ 处的装饰。它由 $\vartheta$ 与 $j$ 典范地确定。微妙之处在于：当 $\zeta$ 绕 $0$ 一圈时，这个小截面一般带有非平凡的单值；由此可见，它实际依赖于 $\vartheta \in \IR$，而不只是 $\vartheta \in \IR / 2 \pi \IZ$。

\item $\gamma^\vartheta_E$ 与 $\CX^\vartheta_{\gamma}$ 的定义与只有正则奇点的情形完全相同。当 $E$ 是边界边时，我们不定义这个圈。

\item 我们把不规则奇点 $\CP_*$ 处的 pop 定义为该处装饰的 $\IZ$-挠子上按 $1$ 的作用。为确定约定：若 $\dots, Q_j, Q_{j+1}, \dots$ 按顺时针排列，而装饰给它们配上截面 $\dots, s_n, s_{n+1}, \dots$，则 pop 之后配上的是 $\dots, s_{n-1}, s_{n}, \dots$。若把 $\vartheta$ 换成 $\vartheta + \pi$，则 $T_\WKB(\vartheta)$ 在每个不规则奇点处的装饰都经历一次 pop。

\item 可能出现这样的情况：一串 flip 产生的新三角剖分，与原来的只差某个边界圆旋转 $2\pi/L$，等价地说，只差 $Q_i$ 的一个循环置换。这与在不规则顶点处作一次 pop 的效果完全相同。这反映了装饰三角剖分之群胚中态射之间的一个关系。

\end{enumerate}

\section{线性 $SU(2)$ 箭图的标度极限：单个不规则奇点的情形}\label{section:Superconformal}

% TODO 图注中「第 9 节」为硬编码（原文 \protect\ref{section:Superconformal}），合稿时核对
\insfig{3.5in}{quiver-of-twos}{第 9 节所考虑理论的线性箭图。}

本节用一些格外简单的理论来说明第 5–8 节% TODO 合稿改 \ref{sec:fg-coordinates}–\ref{sec:irregular}
的某些考虑。这些理论是 $n$ 个 $SU(2)$ 规范群的线性箭图的某些标度极限，其中首、末两个规范群各带两个基本超多重态，如图 43 所示% TODO 合稿改 \ref{fig:quiver-of-twos}
。相应的 Hitchin 系统规范群为 $SU(2)$，并且在标度极限之后，在 $C = \IC\IP^1$ 上只有一个奇点，而且是不规则奇点。

\subsection{线性 $SU(2)$ 箭图及其参数空间}

这个理论有 $3n+3$ 个物理参数：紫外拉氏量中含 $n$ 个规范耦合、$n-1$ 个双基本表示质量与 $4$ 个基本表示质量；而库仑分支由 $n$ 个真空期望值 $\langle {\Tr}(\Phi^{(\alpha)})^2 \rangle$ 参数化。

回忆第 3 节% TODO 合稿改 \ref{sec:physics}
：相应的 Seiberg–Witten 曲线是 $F(t, v) = 0$，其中
\begin{equation}\label{eq:ch11-Ftv}
F(t,v) = \sum_{\alpha=0}^{n+1} q_\alpha (v) t^{n+1-\alpha} = p_0(t) v^2 + p_1(t) v + p_2(t).
\end{equation}
此外我们已经看到，可以（带有某些冗余地）参数化为
\begin{equation}\label{eq:ch11-qalpha}
q_\alpha(v) = c_\alpha (v^2 - \mu_\alpha v - u_\alpha).
\end{equation}
在弱耦合区域，耦合由 $c_\alpha$ 确定，质量由 $\mu_\alpha$ 与 $u_\alpha$ 确定，而库仑分支由 $u_\alpha$ 参数化。最后，正如我们在 (3.77) 处看到的% TODO 跨块引用
，把质心自由度分离出去之后，Seiberg–Witten 曲线就成为某个 $SU(2)$ Hitchin 系统的谱曲线，其中
\begin{equation}\label{eq:ch11-lambda2}
\lambda^2 = \half \frac{p_1(t)^2 - 4 p_2(t) p_0(t) }{(t p_0(t))^2} dt^2.
\end{equation}
这个二次微分在 $t = 0$ 与 $t = \infty$ 处有二阶极点，在 $p_0(t)$ 的 $n+1$ 个零点处也有二阶极点。

以上描述多少掩盖了问题的 S-对偶性质 \cite{Witten:1997sc,Gaiotto:2009we}。物理规范耦合其实只依赖于带 $n+3$ 个标号点的球面之模空间中的一个点。而且，双基本表示质量参数与基本表示质量参数之间并没有 S-对偶不变的区别。这些事实提示：我们应当更一视同仁地对待这 $n+3$ 个奇点与质量参数。办法是在描述中引入一些冗余，即从 $[t:1] \in \IC\IP^1$ 到新坐标 $z$ 作一般的分式线性变换。作这样的变换之后我们得到（多项式下标表示其次数）
\begin{equation}\label{eq:quad-diff-1}
\lambda^2 = \frac{Q_{2n+2}(z)}{(D_{n+3}(z))^2} (dz)^2.
\end{equation}
扣掉 $SL(2,\IC)$ 作用的 $3$ 个参数与 $Q$、$D$ 同时重标度的 $1$ 个参数之后，仍有 $3n + 7 - 4 = 3n+3$ 个物理参数。

如已提到，这些物理参数的地位并不平等。耦合与质量指定紫外理论，而库仑分支参数指定真空。为凸显这一区别，把 $\lambda^2$ 稍作不同的参数化是有用的。

$n+3$ 个质量参数可以刻画为 $\lambda^2$ 中 $n+3$ 个极点的留数。（我们不考虑其中某些极点相撞的情形——那对应物理理论的强耦合奇点。）若这些极点位于 $z_a$，$a = 1, \dots, n+3$，并假定所有 $z_a \neq \infty$，则可以把 $\lambda^2$ 写成如下形式：
\begin{equation}\label{eq:part-frac}
\lambda^2 = \sum_{a=1}^{n+3} \left(
\frac{m_a^2}{(z-z_a)^2} + \frac{c_a}{z-z_a}\right) dz^{2}.
\end{equation}
要求 $\lambda^2$ 在 $z = \infty$ 处正则，给出对 $c_a$ 的三个条件：
\begin{equation}\label{eq:c-conds}
\begin{aligned}
\sum_{a=1}^{n+3} c_a & = 0, \\
\sum_{a=1}^{n+3} z_a c_a & =   - \sum_{a=1}^{n+3} m_a^2, \\
\sum_{a=1}^{n+3} z_a^2 c_a & =   - 2\sum_{a=1}^{n+3} m_a^2 z_a.
\end{aligned}
\end{equation}
库仑分支 $\CB$ 就是满足 \eqref{eq:c-conds} 的 $c_a$ 构成的空间。由于 \eqref{eq:c-conds} 是关于 $c_a$ 的非齐次线性方程，它是一个 $n$ 维仿射空间。

（另一种看法是：一旦通过 $z_a$ 指定了耦合，带任意质量与真空的理论空间就是一个线性空间——多项式 $Q_{2n+2}(z)$ 的空间——它也可以看作 $H^0(\IC \IP^1, K^{\otimes 2} \otimes \CO(2\CP_1 + \cdots + 2\CP_{n+3}) )$。固定质量于是固定这个空间的一个仿射子空间。）

\subsection{标度极限}

对一般的质量，库仑分支的奇点出现在 $Q_{2n+2}(z)$ 的两个或多个根重合之时。在这些点上库仑分支的度量奇异，绕该轨迹的涨落具有无穷大的作用量。原因当然是：在这个轨迹上至少有一个 BPS 态变成无质量。$Q_{2n+2}(z)$ 的零点落在判别式轨迹 $\CD$ 上，$\CD$ 带有分层 $\CD_k$，其中恰有 $k$ 个零点重合。下面我们集中看一个分层 $\CD_N$。在该轨迹附近，若干个 BPS 态变轻，我们想关注这些态的低能物理。

相应地，我们考虑如下的标度极限。取一点 $\CP \in C$ 及局部坐标 $z$，使 $z(\CP) = 0$。$Q$ 的空间中的有关邻域可以用 $Q_{2n+2}(z) = \prod_{i=1}^N (z - \epsilon \theta_i) \tilde Q_{2n+2-N}(z)$ 参数化，其中 $\tilde Q_{2n+2-N}(0) \not= 0$。当 $\epsilon \to 0$ 时，$Q$ 的 $N$ 个零点在 $z = 0$ 处相撞。与连接其中两个相撞零点的弦相应的典型 BPS 态，其质量为 $\epsilon^{(N+2)/2}$ 阶。于是我们考虑这一能标及以下的物理。我们还用 $z = \epsilon \tilde z$ 定义黎曼面 $C$（即 M5-膜世界体）上的一个标度区域；我们只考虑定域在这个区域内的涨落。对 Seiberg–Witten 微分 $\lambda$ 与 $\tilde z$ 作适当重标度之后，我们于是考虑一个理论，其中
\begin{equation}\label{eq:PolyDegen}
\lambda^2 = P_N(z) dz^2,
\end{equation}
$P_N$ 是 $N$ 次首一多项式，而我们这片 M5-膜世界体上的标度区块已被放大成整个复平面。

原理论的许多形变在标度极限下变得不可归一化；不过有些形变仍然可归一化。它们就是多项式形变 $P_N \to P_N + \delta P_N$，满足
\begin{equation}\label{eq:ch11-normaliz}
\int_{\IC} \vert \delta \lambda \vert^2 = \int_{\IC} \biggl\vert
\frac{\delta P_N(z) }{2 \sqrt{P_N(z)}} \biggr\vert^2 d^2 z < \infty.
\end{equation}
由于 $P_N(z)$ 的零点是单零点，发散只可能来自 $z \to \infty$，因此可归一化条件为
\begin{equation}\label{eq:ch11-deltaP}
\delta P_N = \sum_{i < \half(N-2)} \delta p_i z^i.
\end{equation}

对这个可归一化条件还有另一种有用的看法，即固定 $\sqrt{P_N(z)}$ 在 $z = \infty$ 附近展开的奇异部分：
\begin{equation}\label{eq:AsympOfP}
\sqrt{P_N(z)} = \Delta(z) + o(z^{-1})
\end{equation}
其中 $\Delta(z)$ 固定（由我们多半已经丢弃的紫外理论决定）。

$N$ 为偶与 $N$ 为奇之间有一个重要差别。当 $N$ 为偶时，$\Delta(z)$ 是整数幂的展开，以一个单极点收尾；我们把它的留数记作 $m$。变动 $m$ 是「对数可归一化」形变，对应 $i = \half(N-2)$。稍后我们从 Hitchin 系统过渡到平坦联络时就会看到，$m$ 指定了 $\infty$ 处的形式单值。当 $N$ 为奇数时，不存在这样的对数可归一化形变。

标度极限之后，库仑分支上仍然存在奇点轨迹，使某些 $\theta_i = \theta_j$，即 $P_N(z)$ 有多重零点。这些轨迹是 $\CD$ 的高维分层与 $\CD_N$ 周围标度区域的交。在这些轨迹上某些 BPS 态变成无质量。而且，在有些情形下变成无质量的 BPS 态彼此非定域，从而导致 \cite{Argyres:1995jj,Argyres:1995xn} 中所研究的那类理论。

要到达最极端的情形，我们可以调节不可归一化参数，使得所考虑的是 $\lambda^2 = z^N (dz)^2$ 的可归一化形变。模空间中这一点的深红外物理由一个超共形场论描述 \cite{Argyres:1995jj,Argyres:1995xn}；并且可以用上面找出的可归一化形变对此理论作微扰，从而定义一个具有有限 BPS 谱的紫外完备量子场论。于是我们研究的就是这些理论的 Seiberg–Witten 曲线与 BPS 态。从由线性箭图指定的原紫外理论的观点看，我们关注的是模空间某个区域内的一个低能子部门，在那里若干 BPS 态正变得参数式地轻。如下文将看到的，这类态的数目被 $\half N(N-1)$ 界住 \cite{Shapere:1999xr}。

\subsection{Hitchin 系统}\label{subsec:HitchinForAD}

上一小节我们描述了 Hitchin 系统谱曲线的一种退化。为使讨论完整，现在说明 Hitchin 系统相应的退化。新的 Hitchin 系统定义在 $\IC\IP^1$ 上，在 $z=\infty$ 处只有一个不规则奇点。具体地说，当 $N$ 为偶时我们有
\begin{equation}\label{eq:ch11-A0}
A_0 = \begin{pmatrix} -m^{(3)} & 0 \\ 0 & m^{(3)}
\end{pmatrix} \left( \frac{dz}{z} - \frac{d \bar z}{\bar z} \right)
\end{equation}
以及
\begin{equation}\label{eq:ch11-varphi0}
\varphi_0 = \begin{pmatrix} \Delta(z) & 0 \\  0 & - \Delta(z)
\end{pmatrix},
\end{equation}
其中 $\Delta(z)$ 如 \eqref{eq:AsympOfP} 所定义。由于 $\Delta(z) \sim z^{N/2} + \cdots + \frac{m}{z}$，逆时针绕行下的形式单值为
\begin{equation}\label{eq:MonodromyAtInfinity}
\exp \left(2\pi i \left[ \frac{mR}{\zeta} - 2 m^{(3)} - \bar m R
\zeta \right]\sigma^3 \right).
\end{equation}
其特征值我们记作 $\mu^{\pm 1}$。

当 $N$ 为奇数时，没有与质量参数 $m, m^{(3)}, \bar m$ 类似的东西。取而代之，我们有 (3.115) 的一个推广% TODO 跨块引用
：
\begin{align}\label{eq:NoddBC}
\varphi_0 & = \Delta(z) \begin{pmatrix} 0 & (\bar z/z)^{1/4} \\
(z/\bar z)^{1/4} & 0 \end{pmatrix}, \\
A_0 & = \frac{1}{8} \sigma^3
\left(\frac{dz}{z} - \frac{d\bar z}{\bar z}\right).
\end{align}

\begin{yizhu}
$\IZ$-挠子（$\IZ$-torsor）指以 $\IZ$ 为结构群的主齐性空间：它本身没有自然的原点，但任意两个元素之差是一个整数。装饰的选取构成 $\IZ$-挠子，意味着「选定一个装饰」与「选定一个整数偏移」是等价的。
\end{yizhu}
